\documentclass{article}
\usepackage[preprint]{neurips_2025}

\usepackage[utf8]{inputenc}
\usepackage[T1]{fontenc}
\usepackage{hyperref}
\usepackage{url}
\usepackage{booktabs}
\usepackage{amsfonts, amsmath, amssymb}
\usepackage{nicefrac}
\usepackage{microtype}
\usepackage{xcolor}
\usepackage{graphicx}
\usepackage{subcaption}
\usepackage{natbib}

\title{Update Objects for Language Model Training: Replaying Training Progress Instead of Tokens}

\author{
  Julio Jose Lena \\
  Independent Researcher \\
  \texttt{juliolena25@gmail.com} \\
  \url{https://orcid.org/0009-0001-0817-6591}
}

\begin{document}

\maketitle

\begin{abstract}
Large language model training still treats text tokens as the main reusable
unit of learning. If a training interval has already transformed parameters
from $\theta_t$ to $\theta_{t+k}$, then the interval also emits a causal object:
the state transition that the optimizer actually bought. This paper introduces
\emph{update objects}, compact artifacts that replay useful training movement
without replaying all source tokens. I study two concrete objects in small
conventional decoder-only Transformers: sparse checkpoint delta replay objects
and optimizer transaction capsules. All target arms start from the same
$\theta_t$ and are compared against raw token replay, gradient-alignment
selected replay, checkpoint distillation, no-source continuation, warm-start
upper-bound controls, and shuffled-object ablations. Sparse checkpoint deltas
are the strongest result: across two model widths, two text shards, compression
sweeps, pessimistic build-cost accounting, and a five-seed d96 paper-core gate,
they improve durable held-out cross-entropy progress per paid unit by
$2.45$--$2.90\times$ over the best non-object replay control. A more selective
secondary object also survives: optimizer transaction capsules beat raw replay
by $1.78\times$ in the same d96 five-seed paper-core gate, while mean-gradient,
layerwise-gradient, alignment-weighted, and sign-consensus capsules fail or
remain weak. These results do not claim production-scale speedup. They support
a narrower thesis: previous training runs can produce reusable transition
objects that are empirically different from data selection, distillation,
checkpoint warm starts, or parameter-efficient adaptation.
\end{abstract}

\section{Introduction}

Modern language models are trained by repeatedly presenting text to an
optimizer, measuring loss, and changing parameters. Once the step is complete,
the tokens, gradients, optimizer state, and parameter movement are usually
discarded or reduced to logs. Later work that needs the same capability must
typically replay the text, select a smaller subset of text, distill from a
teacher, warm-start from a checkpoint, or fine-tune a low-rank adapter. Those
methods reuse examples, functions, or weights. They rarely treat the training
transition itself as reusable data.

This paper asks a different question: after a corpus block has already moved a
model from $\theta_t$ to $\theta_{t+k}$, can the useful part of that movement be
stored as a compact object and replayed into another compatible training state?
The motivating unit is not a token and not a teacher logit. It is an update
object: an artifact whose payload is a priced, causal summary of what the
training interval did to the model.

The distinction is important. Data selection methods choose which examples to
train on; update objects attempt to preserve the effect of a completed training
interval. Dataset distillation synthesizes examples that induce useful future
updates; update objects store the update-side artifact directly. Function
preserving growth and model merging reuse trained functions or parameter
states; update objects are evaluated as replayable training progress under an
explicit cost denominator. This boundary makes the proposal testable: an update
object should buy more durable held-out cross-entropy reduction per paid unit
than replaying the original tokens or using nearby reuse baselines.

The practical motivation is not that token replay is obsolete. Tokens remain
the ground truth interface for language. The claim is narrower: modern training
runs already produce expensive intermediate evidence about how a model changes
when a block of data is consumed. If that evidence can be stored as a compact
transition object, then failed runs, pilot runs, refresh runs, and corpus
passes may become a source of reusable training data rather than only sunk
cost.

I make four contributions.
\begin{enumerate}
  \item I define \emph{update objects} as priced, replayable summaries of a
  training-state transition, and separate them from examples, logits, adapters,
  checkpoint warm starts, and ordinary model diffs.
  \item I introduce two concrete objects for conventional Transformer gates:
  sparse checkpoint delta replay objects and optimizer transaction capsules.
  \item I define a durable cross-entropy progress score that charges object
  build, application, replay, selection, distillation, and continuation costs in
  token-backward-equivalent units.
  \item I report falsifier-style gates in small decoder-only Transformers,
  showing that checkpoint deltas are consistently strong and that only the
  optimizer-transaction version of capsules survives mature controls.
\end{enumerate}

The results are intentionally modest in scale but strict in control design.
Every arm begins from the same $\theta_t$. Object arms are compared to raw
token replay, gradient-alignment selected replay, checkpoint distillation,
no-source continuation, warm-start upper bounds, and shuffled-object
ablations. The most important run is a d96, five-seed paper-core gate:
checkpoint delta objects score $3.624947$ durable CE progress per paid unit,
optimizer transaction capsules score $2.231399$, and raw replay scores
$1.250833$. The corresponding ratios are $2.898\times$ and $1.784\times$
against raw replay.

\section{Related Work}

\paragraph{Data selection and influence.}
LESS selects influential instruction-tuning examples by building reusable
low-rank gradient features and ranking examples by similarity to a target
capability \citep{xia2024less}. This paper uses a similar intuition that
gradient-side information is valuable, but the object boundary is different:
the selected unit is not an example. The replayed unit is a parameter-space
transition object emitted by a source interval.

\paragraph{Dataset distillation and trajectory matching.}
Dataset distillation compresses a dataset into synthetic examples that train a
model efficiently \citep{wang2018datasetdistillation}. Trajectory matching
pushes distilled data so that synthetic training follows expert training
trajectories \citep{cazenavette2022trajectory}. Update objects move one step
closer to the trajectory itself. Rather than synthesizing inputs that induce a
trajectory, they store sparse pieces of the trajectory or optimizer
transaction and test whether those pieces can be replayed directly.

\paragraph{Function-preserving growth and parameter reuse.}
Net2Net transfers a function into wider or deeper networks
\citep{chen2015net2net}, and LiGO learns growth operators for pretrained
Transformers \citep{wang2023ligo}. LoRA stores task adaptation in low-rank
parameter updates \citep{hu2022lora}. Task arithmetic defines task vectors as
differences between pretrained and fine-tuned weights and uses arithmetic in
weight space to edit model behavior \citep{ilharco2023taskarithmetic}.
TIES-Merging trims small deltas, resolves sign conflicts, and merges only
parameters aligned with an elected sign \citep{yadav2023ties}. DARE drops and
rescales delta parameters before merging homologous language models
\citep{yu2024supermario}. These methods are the closest prior-art pressure for
the checkpoint-delta result. Update objects are not claimed to replace them,
and the paper should not be read as saying that sparse deltas are new by
themselves. The distinction is the experimental contract: the object is scored
as replayable paid training progress from a source interval, against token
replay, selected replay, distillation, continuation, warm-start collapse, and
shuffled-object controls.

\section{Update Objects}

\subsection{Training Intervals}

Let $\theta_t$ be a model state and let $B_{t:t+k-1}$ be a sequence of source
minibatches. A training rule $U$ maps
\begin{equation}
  \theta_{i+1} = U(\theta_i, B_i)
\end{equation}
for $i=t,\ldots,t+k-1$, producing a terminal state $\theta_{t+k}$. The usual
record of this interval is the source data and perhaps a checkpoint. An update
object treats the interval as an object-producing computation:
\begin{equation}
  O_{t,k} = E(\theta_t, \theta_{t+k}, B_{t:t+k-1}, U, M),
\end{equation}
where $E$ is an encoder and $M$ denotes optional measured side information such
as gradients, optimizer steps, or coordinate-wise movement. The object is
useful only if a replay operator $R$ can apply it to a compatible state:
\begin{equation}
  \widehat{\theta}_{t+k} = R(\theta_t, O_{t,k}).
\end{equation}
This paper does not require exact reconstruction of $\theta_{t+k}$. The target
is cheaper durable progress on held-out cross-entropy after replay and a short
continuation interval.

\subsection{Durable Progress Per Paid Unit}

Let $D_{\mathrm{eval}}$ be a held-out evaluation split and let
$\mathrm{CE}(\theta;D_{\mathrm{eval}})$ be mean next-token cross-entropy. The
initial held-out loss is
\begin{equation}
  L_0 = \mathrm{CE}(\theta_t;D_{\mathrm{eval}}).
\end{equation}
For any arm $a$, the gate first applies the arm-specific replay operation,
then trains for the same continuation budget, producing $\theta_a^{+}$. The
durable progress is
\begin{equation}
  \Delta_{\mathrm{dur}}(a) =
  L_0 - \mathrm{CE}(\theta_a^{+};D_{\mathrm{eval}}).
\end{equation}
The paid score divides this durable progress by an accounting denominator:
\begin{equation}
  S(a) = \frac{\Delta_{\mathrm{dur}}(a)}
  {C_{\mathrm{replay}}(a)+C_{\mathrm{build}}(a)+C_{\mathrm{apply}}(a)
  +C_{\mathrm{select}}(a)+C_{\mathrm{distill}}(a)}.
  \label{eq:paid_score}
\end{equation}
The costs in this paper are token-backward-equivalent proxy units. They are not
claimed to be hardware FLOPs. The purpose is to prevent an object from winning
by hiding its preprocessing or teacher cost. A later systems paper would need
wall-clock or profiler-derived FLOP accounting; the present paper only tests
whether the object mechanism survives fair first pricing.

\subsection{Admissible Controls}

The control set is part of the definition. An update object is not meaningful
if it only beats doing nothing. It should beat or clearly complement:
\begin{itemize}
  \item \textbf{raw replay}, which trains on the original source batches;
  \item \textbf{selected replay}, which pays for gradient-alignment selection
  and trains on the selected source batches;
  \item \textbf{distillation}, which trains from $\theta_t$ against
  $\theta_{t+k}$ logits on the same source batches;
  \item \textbf{no-source continuation}, which spends only the continuation
  budget without source replay;
  \item \textbf{warm-start checkpoint}, an upper-bound collapse control that
  starts from $\theta_{t+k}$ but is charged a larger source-build cost;
  \item \textbf{shuffled ablations}, which preserve object density and
  magnitude distribution while destroying coordinate identity.
\end{itemize}

These controls make the paper falsifiable. A method that loses to raw replay is
not a useful replay object. A method that loses to its shuffled ablation is not
using the claimed coordinate information. A method that only wins by ignoring
build or selection cost is not an update-object gain.

\section{Two Concrete Update Objects}

\subsection{Sparse Checkpoint Delta Replay Objects}

The first object is the simplest transition summary:
\begin{equation}
  \Delta_{t,k} = \theta_{t+k} - \theta_t.
\end{equation}
A full checkpoint delta is too close to warm-starting the target from a source
checkpoint, so the gate uses a sparse transition object. Let
$\mathrm{TopK}_\rho(v)$ keep the largest $\rho$ fraction of coordinates by
absolute value and set all other coordinates to zero. The checkpoint delta
object is
\begin{equation}
  O^{\Delta}_{t,k} = \mathrm{TopK}_\rho(\Delta_{t,k}),
  \qquad
  R_{\Delta}(\theta_t,O^{\Delta}_{t,k}) =
  \theta_t + \alpha_{\Delta} O^{\Delta}_{t,k}.
  \label{eq:checkpoint_delta_object}
\end{equation}
In the gates below, $\alpha_{\Delta}=1$. The object is therefore a sparse,
coordinate-preserving replay of the state movement bought by the source
interval. Its shuffled ablation keeps the same values and density but permutes
coordinates. This is the key sanity check: if the object only works because it
adds noise or changes norm, the shuffled version should also work.

This object is intentionally vulnerable to prior-art pressure. It resembles
checkpoint compression, model-diff transfer, task arithmetic, DARE-style delta
sparsification, and sparse model merging. The paper's escape boundary is not
``a sparse delta exists.'' The boundary is that the delta is scored as a
replayable training-transition object under durable CE progress per paid unit,
against token replay and distillation controls.

\subsection{Optimizer Transaction Capsules}

The second object records optimizer-side movement rather than the endpoint
difference. Let the source interval produce step transactions
\begin{equation}
  u_i = \theta_{i+1} - \theta_i.
\end{equation}
The signed transaction over the interval is
\begin{equation}
  T_{t,k} = \sum_{i=t}^{t+k-1} u_i,
\end{equation}
and the coordinate activity score is
\begin{equation}
  A_{t,k}[j] = \sum_{i=t}^{t+k-1} |u_i[j]|.
\end{equation}
The optimizer transaction capsule keeps the signed transaction values on the
coordinates with largest cumulative movement:
\begin{equation}
  O^{\mathrm{opt}}_{t,k}[j] =
  \begin{cases}
    T_{t,k}[j] & \text{if } j \in \mathrm{TopK}_{\rho}(A_{t,k}), \\
    0 & \text{otherwise.}
  \end{cases}
  \label{eq:optimizer_capsule}
\end{equation}
The replay operator is direct addition,
\begin{equation}
  R_{\mathrm{opt}}(\theta_t,O^{\mathrm{opt}}_{t,k})
  = \theta_t + \alpha_{\mathrm{opt}} O^{\mathrm{opt}}_{t,k},
\end{equation}
with $\alpha_{\mathrm{opt}}=1$ in the reported gates.

The activity score matters. A mean gradient capsule stores a sparse average
gradient and applies a single learning-rate-scaled update. That object was too
weak in the gates. The surviving capsule is different: it stores the sparse
sum of actual optimizer transactions and selects coordinates by cumulative
movement. In other words, the object preserves what the optimizer repeatedly
used during the source interval, not merely what the average gradient suggested
at the beginning.

\subsection{Why the Objects Are Not Equivalent}

For plain gradient descent with fixed learning rate, no momentum, no adaptive
state, and no nonlinearity from changing batches, the sum of optimizer
transactions can collapse toward a checkpoint delta. In realistic training,
the objects separate. The checkpoint delta is an endpoint summary. The
optimizer capsule is a transaction summary with its own activity criterion.
The empirical gates reflect this separation: checkpoint deltas are stronger,
but optimizer transaction capsules survive where gradient-only capsules do not.

\section{Experimental Protocol}

\subsection{Models and Data}

The experiments use small conventional decoder-only Transformers implemented in
PyTorch. The gates are deliberately not ternary-SSM gates; the purpose is to
test whether the update-object idea has a conventional-language-model form. I
use model widths d48, d64, and d96, AdamW optimization, byte-token training
shards derived from Wikipedia-like text, and held-out next-token
cross-entropy. The reported tables use two shard sizes: a smaller
\texttt{wiki\_1mb} shard for early falsifiers and a larger \texttt{wiki\_5mb}
shard for the more stable gates.

Each run has the same high-level structure. First, a source interval trains a
model from $\theta_t$ to $\theta_{t+k}$ and emits candidate update objects.
Second, every replay arm is reset to the same $\theta_t$. Third, each arm
applies its replay operation or source replay. Fourth, each arm receives the
same continuation budget. The final score is computed from held-out CE after
this continuation phase. This design avoids the most common false positive:
giving an object arm a hidden warm start or a different continuation budget.

\subsection{Arms}

The core arms are:
\begin{itemize}
  \item \textbf{raw replay}: replay the source batches with ordinary AdamW;
  \item \textbf{selected replay}: replay a selected subset of source batches
  chosen by gradient alignment, with selection cost charged;
  \item \textbf{checkpoint distillation}: train from $\theta_t$ against
  $\theta_{t+k}$ teacher logits on the same source batches;
  \item \textbf{no-source continuation}: spend only the continuation budget;
  \item \textbf{checkpoint delta object}: apply
  Eq.~\ref{eq:checkpoint_delta_object};
  \item \textbf{optimizer transaction capsule}: apply
  Eq.~\ref{eq:optimizer_capsule};
  \item \textbf{gradient capsule variants}: mean-gradient, layerwise-gradient,
  alignment-weighted, and sign-consensus capsules;
  \item \textbf{shuffled ablations}: density- and magnitude-matched shuffled
  versions of the object arms;
  \item \textbf{warm-start checkpoint}: initialize from $\theta_{t+k}$ and
  charge a larger source-build cost.
\end{itemize}

Warm-start is not treated as a deployable object in the main claim. It is a
collapse control: if an object merely approximates starting from the terminal
checkpoint, then warm-start should dominate when the source-build cost is not
counted. The stricter question is whether a compact object can beat replay
after its own construction and application costs are charged.

\subsection{Cost Accounting}

The raw replay denominator is the number of source replay token-backward units
plus the shared continuation units. Object arms pay a build cost, an
application cost, and the same continuation cost. Selection and distillation
arms pay their own selection or teacher costs. Several gates deliberately vary
the checkpoint-delta build cost to test whether the result survives pessimistic
accounting. The resulting score is therefore a \emph{paid-unit} gate, not a
claim about optimized wall-clock training time.

This distinction is central to the paper. A production system could implement
objects inefficiently and lose wall-clock time; conversely, a kernel-level
implementation might make object application cheaper than the proxy used here.
The empirical claim is only that, under explicit first pricing, some transition
objects buy more durable held-out CE progress than replaying the source tokens.

\section{Sparse Checkpoint Delta Results}

Table~\ref{tab:checkpoint_delta_main} reports the main checkpoint-delta gates.
The object wins in every listed condition. The first three runs use three
seeds. The paper-lite follow-up uses five seeds on \texttt{wiki\_5mb} with
d64, and the checkpoint delta object wins all five seeds against the best
non-object replay/control.

\begin{table}[t]
  \centering
  \caption{Sparse checkpoint delta replay objects. Scores are durable held-out
  CE progress per paid unit.}
  \label{tab:checkpoint_delta_main}
  \begin{tabular}{llrrrrrr}
    \toprule
    Run & Data & d & Tokens & Ratio & Delta & Raw & Distill \\
    \midrule
    d48 main & \texttt{wiki\_1mb} & 48 & 60k & 2.487 & 0.600784 & 0.241524 & 0.199657 \\
    d64 main & \texttt{wiki\_1mb} & 64 & 80k & 2.456 & 1.206092 & 0.491047 & 0.385313 \\
    d64 larger shard & \texttt{wiki\_5mb} & 64 & 150k & 2.454 & 1.215255 & 0.495229 & 0.390310 \\
    d64 paper-lite & \texttt{wiki\_5mb} & 64 & 180k & 2.468 & 1.165835 & 0.472301 & 0.368169 \\
    \bottomrule
  \end{tabular}
\end{table}

The result is not just a win over an unusually weak baseline. In the same
runs, gradient-alignment selected replay is close to raw replay but does not
beat the checkpoint delta object, and checkpoint distillation is weaker still.
The shuffled-delta ablation is also lower than the true coordinate-preserving
delta. In the paper-lite run, for example, the shuffled-delta score is
$0.400597$ while the true delta score is $1.165835$. This gap matters because
it shows that coordinate identity is doing work; the object is not only adding
a sparse perturbation with the right magnitude distribution.

\subsection{Compression and Build-Cost Sensitivity}

Table~\ref{tab:checkpoint_delta_sensitivity} tests whether the checkpoint
delta object survives stronger compression or pessimistic build accounting.
At density $0.005$, the ratio falls but remains positive at $1.871\times$.
At density $0.05$, the ratio rises to $2.638\times$. When the object build
cost is raised to ten paid units, the ratio drops to $1.379\times$ but remains
above raw replay. This is important for paper honesty: the object is not free,
and a sufficiently expensive implementation could erase the gain. The current
claim is that the mechanism survives first-order pricing pressure.

\begin{table}[t]
  \centering
  \caption{Checkpoint-delta sensitivity gates.}
  \label{tab:checkpoint_delta_sensitivity}
  \begin{tabular}{llr}
    \toprule
    Gate & Change & Object/control ratio \\
    \midrule
    Compression 0.005 & stronger sparsity & 1.871 \\
    Compression 0.05 & more transition mass & 2.638 \\
    Build cost 10 & pessimistic object build cost & 1.379 \\
    \bottomrule
  \end{tabular}
\end{table}

The stronger d96 paper-core gate appears in the combined results below because
it also includes the mature capsule family. Its checkpoint-delta result is the
strongest single measurement in the paper: $3.624947$ paid progress versus
$1.250833$ for raw replay, a $2.898\times$ ratio over the best non-object
control.

\section{Optimizer Transaction Capsule Results}

The initial capsule result was weak because the first capsule was effectively a
sparse average-gradient object. That object did not beat raw replay. The mature
gate therefore separates gradient-only capsules from optimizer transaction
capsules. Table~\ref{tab:capsule_mature} summarizes the surviving object and
the failure modes.

\begin{table}[t]
  \centering
  \caption{Mature optimizer transaction capsule gates. Scores are durable
  held-out CE progress per paid unit.}
  \label{tab:capsule_mature}
  \begin{tabular}{llrrrrr}
    \toprule
    Run & Model & Density & Capsule & Raw & Shuffled & Ratio \\
    \midrule
    source24 & d64 & 0.02 & 0.506960 & 0.472301 & 0.380432 & 1.073 \\
    source48 & d64 & 0.02 & 0.635740 & 0.489593 & 0.398688 & 1.298 \\
    source48 & d64 & 0.05 & 0.904204 & 0.489593 & 0.399149 & 1.847 \\
    source48 & d64 & 0.005 & 0.483575 & 0.489593 & 0.398317 & 0.988 \\
    build-cost 10 & d64 & 0.05 & 0.561091 & 0.489593 & 0.247686 & 1.146 \\
    d96 scale check & d96 & 0.05 & 2.267360 & 1.261991 & 0.971832 & 1.797 \\
    d96 paper-core & d96 & 0.05 & 2.231399 & 1.250833 & 0.987417 & 1.784 \\
    \bottomrule
  \end{tabular}
\end{table}

The capsule is therefore not as strong as the checkpoint delta object, but it
is not dead. It has a clear operating region. It needs enough source interval
length, enough density, and direct access to optimizer transactions. It fails
when over-compressed to density $0.005$. It remains positive under a
pessimistic build-cost setting but drops below a strict $1.20\times$ paper
threshold. The d96 paper-core gate is the cleanest result: the capsule scores
$2.231399$, raw replay scores $1.250833$, and the shuffled optimizer capsule
scores $0.987417$.

\subsection{Gradient-Only Capsules Are Not Enough}

The mature gate also tested mean-gradient, layerwise-gradient,
alignment-weighted, and sign-consensus capsules. These are useful negative
controls. In the d96 paper-core gate, their shuffled and unshuffled scores sit
near the same low band, around $0.987$ paid progress for the shuffled variants.
They do not establish a viable update object. This narrows the paper's claim:
``capsules work'' would be too broad, but ``optimizer transaction capsules can
survive when the object records actual optimizer movement'' is supported by the
gates.

This distinction is mathematically plausible. A gradient-only object records a
local direction estimate. An optimizer transaction object records the realized
path after AdamW scaling, batch order, adaptive state, and repeated coordinate
use. If the goal is to replay a training interval rather than predict a single
step, the realized transaction is the more faithful object.

\section{Combined Paper-Core Gate}

The final d96 paper-core gate uses five seeds, source/replay lengths $48/12$,
and density $0.05$. It is the most important table because it tests the
strongest checkpoint delta and the mature capsule in the same conventional
Transformer setting.

\begin{table}[t]
  \centering
  \caption{d96 five-seed paper-core gate.}
  \label{tab:paper_core}
  \begin{tabular}{lrr}
    \toprule
    Arm & Paid progress score & Ratio to raw replay \\
    \midrule
    Checkpoint delta object & 3.624947 & 2.898 \\
    Optimizer transaction capsule & 2.231399 & 1.784 \\
    Raw replay & 1.250833 & 1.000 \\
    Shuffled delta ablation & 1.039598 & 0.831 \\
    Shuffled optimizer capsule & 0.987417 & 0.789 \\
    \bottomrule
  \end{tabular}
\end{table}

The combined gate supports two conclusions. First, sparse checkpoint deltas are
the main paper result. They beat raw replay, distillation, selected replay in
earlier gates, and shuffled coordinate controls. Second, optimizer transaction
capsules are a secondary but meaningful result. They are weaker than checkpoint
deltas, but they beat raw replay and their shuffled capsule control in all five
paper-core seeds.

\section{Discussion}

\subsection{Why This Is More Than Checkpoint Compression}

The strongest objection is that a sparse checkpoint delta is merely checkpoint
compression or task-vector reuse. That objection is partially correct and must
not be hidden. The object payload is a sparse parameter delta, and adjacent
literature already treats weight-space deltas as editable, mergeable, and
sparsifiable objects. The new part is the experimental contract around that
payload. It is not evaluated as a storage artifact, a model merge, a task edit,
or a final checkpoint. It is evaluated as a replayable training-transition
object that competes with token replay, selected replay, distillation,
continuation, warm-start collapse controls, and shuffled coordinate ablations
under a paid progress denominator.

This reframing changes the practical question. A lab already stores many
checkpoints and logs from pilot runs, failed runs, curriculum sweeps, data
refreshes, and ablations. If those runs can emit reusable transition objects,
then training history becomes a cache of paid movement. The paper-core result
does not prove that this cache scales to frontier LLMs, but it makes the
question concrete enough to test.

\subsection{Why Capsules Still Matter}

Checkpoint deltas are strong but endpoint-centered. Optimizer transaction
capsules point at a different family of objects: path-centered training
artifacts. A capsule can, in principle, store where the optimizer repeatedly
spent movement even when the endpoint delta hides cancellation, adaptive
rescaling, or path structure. The mature gates show that a naive gradient
capsule is not enough, but an optimizer transaction capsule can survive. That
is a smaller claim than the initial capsule hope, but it is also cleaner.

The next capsule direction should not be another mean-gradient variant. It
should preserve more optimizer semantics: moments, update signs over time,
block-local activity, curvature tags, or compatibility conditions between the
source and replay states. The paper therefore treats capsules as a promising
secondary line rather than the headline result.

\section{Limitations}

The evidence has five important limitations.

First, the experiments are small. The models are conventional decoder-only
Transformers, but they are not production-scale LLMs. The result should be
read as a falsifier-passing mechanism, not as proof of frontier-scale training
acceleration.

Second, the cost denominator is a proxy. Paid units charge source replay,
build, selection, distillation, application, and continuation in
token-backward-equivalent units. They are useful for avoiding free-object
accounting, but they are not profiler-measured FLOPs, memory traffic, or
wall-clock time.

Third, the strongest object is vulnerable to prior-art pressure. Sparse
checkpoint deltas intersect with checkpoint compression, task arithmetic,
model-diff transfer, DARE/TIES-style delta processing, and model merging. The
paper's novelty depends on the update-object contract: durable replay progress
per paid unit against token replay and distillation controls. A larger paper
should include stronger same-budget baselines: magnitude-pruned task vectors,
DARE/TIES-like delta processing, low-rank or SVD deltas, quantized sparse
deltas, and same-storage-budget checkpoint compression.

Fourth, compatibility is only lightly tested. All target arms begin from the
same $\theta_t$ used by the source interval. A future system must define how
far a replay state can drift before an object becomes unsafe or useless. The
next decisive gate should apply the same object to nearby checkpoints, another
seed, or a controlled $\theta_t+\epsilon$ drift state.

Fifth, quality is measured by held-out cross-entropy only. Capability-level
effects, long-horizon behavior, instruction-following quality, and safety
behavior are outside the current gate.

\section{Conclusion}

This paper introduced update objects: compact, priced, replayable summaries of
training-state transitions. In small conventional decoder-only Transformers,
sparse checkpoint delta replay objects consistently beat raw token replay,
gradient-alignment selected replay, checkpoint distillation, continuation
controls, and shuffled coordinate ablations in durable held-out CE progress per
paid unit. The strongest d96 five-seed gate gives a $2.898\times$ ratio over
raw replay. Optimizer transaction capsules are weaker but survive as a
secondary path-centered object, reaching $1.784\times$ over raw replay in the
same gate, while gradient-only capsules remain weak.

The central claim is therefore not that token replay is obsolete, nor that a
production LLM training run will immediately accelerate by the reported
ratios. The central result is that training intervals can emit reusable
transition artifacts whose value is separable from token replay, distillation,
and shuffled perturbations under explicit paid-progress accounting. If this
boundary scales, language model training could treat previous runs as sources
of replayable training progress rather than only as sources of final
checkpoints and logs.

\bibliographystyle{plainnat}
\bibliography{references}

\end{document}
